% ============================================================
%  MTH636 -- Game Theory | Lecture 4
%  Weak Dominance, Auctions & Max-Min Strategies
% ============================================================
\documentclass[12pt, a4paper]{article}

% ============================================================
%  MTH636 -- Game Theory
%  Notes Preamble — Color-Coded, TikZ-Enhanced
% ============================================================

% --- Core packages ---
\usepackage[T1]{fontenc}
\usepackage[utf8]{inputenc}
\usepackage{lmodern}
\usepackage[margin=2.5cm, top=2.8cm, bottom=2.8cm]{geometry}
\usepackage{microtype}

% --- Math ---
\usepackage{amsmath, amssymb, amsthm, mathtools}
\usepackage{bm}        % bold math
\usepackage{cancel}    % strike-through in math

% --- Color & Boxes ---
\usepackage[dvipsnames, svgnames, x11names, table]{xcolor}
\usepackage[most]{tcolorbox}
\tcbuselibrary{breakable, skins, theorems}

% --- TikZ & Diagrams ---
\usepackage{tikz}
\usepackage{tikz-qtree}
\usetikzlibrary{arrows.meta, positioning, shapes.geometric, fit, backgrounds,
                decorations.pathreplacing, calc, matrix, trees}
\usepackage{pgfplots}
\pgfplotsset{compat=1.18}

% --- Tables ---
\usepackage{booktabs}
\usepackage{multirow}
\usepackage{array}
\usepackage{tabularx}
\usepackage{makecell}

% --- Typography & Layout ---
\usepackage{enumitem}
\usepackage{parskip}
\usepackage{fancyhdr}
\usepackage{titlesec}
\usepackage{titletoc}
\usepackage{hyperref}
\usepackage{graphicx}
\usepackage{float}
\usepackage{caption}

% --- Color Palette ---
\definecolor{defblue}{HTML}{1A5276}      % Definition header
\definecolor{defbluebg}{HTML}{D6EAF8}    % Definition background
\definecolor{thmgreen}{HTML}{1E8449}     % Theorem header
\definecolor{thmbg}{HTML}{D5F5E3}        % Theorem background
\definecolor{exorange}{HTML}{B7770D}     % Example header
\definecolor{exbg}{HTML}{FEF9E7}         % Example background
\definecolor{warnred}{HTML}{922B21}      % Warning/Key Insight header
\definecolor{warnbg}{HTML}{FADBD8}       % Warning background
\definecolor{noticegray}{HTML}{2C3E50}   % Note header
\definecolor{noticebg}{HTML}{EBF5FB}     % Note background
\definecolor{taketawayyellow}{HTML}{7D6608} % Key Takeaway header
\definecolor{takeawaybg}{HTML}{FDFEFE}   % Takeaway background
\definecolor{accentpurple}{HTML}{6C3483} % Accent / proof
\definecolor{lightgray}{HTML}{F2F3F4}    % shading

% ============================================================
%  TCOLORBOX ENVIRONMENTS
% ============================================================

% -- Definition (Blue) --
\newtcolorbox{defbox}[2][]{
  enhanced, breakable,
  colback=defbluebg, colframe=defblue, coltitle=white,
  fonttitle=\bfseries\sffamily,
  title={Definition: #2},
  attach boxed title to top left={yshift=-2mm, xshift=4mm},
  boxed title style={colback=defblue, rounded corners},
  arc=4pt, boxrule=1.2pt,
  left=6pt, right=6pt, top=8pt, bottom=6pt,
  #1
}

% -- Theorem / Proposition (Green) --
\newtcolorbox{thmbox}[2][]{
  enhanced, breakable,
  colback=thmbg, colframe=thmgreen, coltitle=white,
  fonttitle=\bfseries\sffamily,
  title={Theorem: #2},
  attach boxed title to top left={yshift=-2mm, xshift=4mm},
  boxed title style={colback=thmgreen, rounded corners},
  arc=4pt, boxrule=1.2pt,
  left=6pt, right=6pt, top=8pt, bottom=6pt,
  #1
}

% -- Proposition (Green variant) --
\newtcolorbox{propbox}[2][]{
  enhanced, breakable,
  colback=thmbg, colframe=thmgreen!70!black, coltitle=white,
  fonttitle=\bfseries\sffamily,
  title={Proposition: #2},
  attach boxed title to top left={yshift=-2mm, xshift=4mm},
  boxed title style={colback=thmgreen!70!black, rounded corners},
  arc=4pt, boxrule=1.2pt,
  left=6pt, right=6pt, top=8pt, bottom=6pt,
  #1
}

% -- Example (Orange/Yellow) --
\newtcolorbox{exbox}[2][]{
  enhanced, breakable,
  colback=exbg, colframe=exorange, coltitle=white,
  fonttitle=\bfseries\sffamily,
  title={Example: #2},
  attach boxed title to top left={yshift=-2mm, xshift=4mm},
  boxed title style={colback=exorange, rounded corners},
  arc=4pt, boxrule=1.2pt,
  left=6pt, right=6pt, top=8pt, bottom=6pt,
  #1
}

% -- Key Insight / Warning (Red) --
\newtcolorbox{warnbox}[2][]{
  enhanced, breakable,
  colback=warnbg, colframe=warnred, coltitle=white,
  fonttitle=\bfseries\sffamily,
  title={Key Insight: #2},
  attach boxed title to top left={yshift=-2mm, xshift=4mm},
  boxed title style={colback=warnred, rounded corners},
  arc=4pt, boxrule=1.2pt,
  left=6pt, right=6pt, top=8pt, bottom=6pt,
  #1
}

% -- Note (Gray-Blue) --
\newtcolorbox{notebox}[1][]{
  enhanced, breakable,
  colback=noticebg, colframe=noticegray!60, coltitle=white,
  fonttitle=\bfseries\sffamily,
  title={Note},
  attach boxed title to top left={yshift=-2mm, xshift=4mm},
  boxed title style={colback=noticegray!60, rounded corners},
  arc=4pt, boxrule=0.8pt,
  left=6pt, right=6pt, top=8pt, bottom=6pt,
  #1
}

% -- Key Takeaway --
\newtcolorbox{takebox}[1][]{
  enhanced, breakable,
  colback=takeawaybg, colframe=taketawayyellow, coltitle=white,
  fonttitle=\bfseries\sffamily,
  title={Key Takeaway},
  attach boxed title to top left={yshift=-2mm, xshift=4mm},
  boxed title style={colback=taketawayyellow, rounded corners},
  arc=4pt, boxrule=1.2pt,
  left=6pt, right=6pt, top=8pt, bottom=6pt,
  #1
}

% -- Intuition --
\newtcolorbox{intubox}[1][]{
  enhanced, breakable,
  colback=accentpurple!8, colframe=accentpurple!60, coltitle=white,
  fonttitle=\bfseries\sffamily,
  title={Intuition},
  attach boxed title to top left={yshift=-2mm, xshift=4mm},
  boxed title style={colback=accentpurple!80, rounded corners},
  arc=4pt, boxrule=0.8pt,
  left=6pt, right=6pt, top=8pt, bottom=6pt,
  #1
}

% ============================================================
%  PAYOFF MATRIX HELPER
%  Usage: inside tabular, use \payoff{a}{b} for (a,b)
% ============================================================
\newcommand{\payoff}[2]{$(#1,\;#2)$}
\newcommand{\nash}[2]{\underline{$#1$},\underline{$#2$}}  % underline NE payoffs
\newcommand{\Nset}{\mathbb{N}}
\newcommand{\Rset}{\mathbb{R}}
\newcommand{\Sset}{\mathcal{S}}
\newcommand{\Ui}{u_i}

% ============================================================
%  SECTION FORMATTING
% ============================================================
\titleformat{\section}
  {\large\bfseries\sffamily\color{defblue}}
  {\thesection}{1em}{}[\titlerule]

\titleformat{\subsection}
  {\normalsize\bfseries\sffamily\color{thmgreen!80!black}}
  {\thesubsection}{1em}{}

\titleformat{\subsubsection}
  {\normalsize\itshape\sffamily\color{exorange!80!black}}
  {\thesubsubsection}{1em}{}

% ============================================================
%  HEADER / FOOTER
% ============================================================
\pagestyle{fancy}
\fancyhf{}
\fancyhead[L]{\sffamily\small\color{defblue} MTH636 -- Game Theory}
\fancyhead[R]{\sffamily\small\color{defblue} \leftmark}
\fancyfoot[C]{\sffamily\small\thepage}
\renewcommand{\headrulewidth}{0.4pt}

% ============================================================
%  HYPERREF SETUP
% ============================================================
\hypersetup{
  colorlinks=true,
  linkcolor=defblue,
  urlcolor=thmgreen,
  citecolor=accentpurple,
  bookmarksnumbered=true,
  pdftitle={MTH636 Game Theory Notes},
  pdfauthor={MTH636 Student Notes}
}

% ============================================================
%  MISC HELPERS
% ============================================================
\newcommand{\BR}{\text{BR}}        % Best Response
\newcommand{\NE}{\text{NE}}        % Nash Equilibrium
\newcommand{\SPNE}{\text{SPNE}}    % Subgame Perfect NE
\newcommand{\Dom}{\text{dom}}      % Dominance
\newcommand{\mixs}[1]{\sigma_{#1}} % Mixed strategy
\newcommand{\pures}[1]{s_{#1}}     % Pure strategy


% ============================================================
\begin{document}

% Title
\begin{center}
  {\LARGE\bfseries\sffamily\color{defblue} MTH636 -- Game Theory}\\[6pt]
  {\Large\sffamily Lecture 4: Weak Dominance, Auctions \& Max-Min Strategies}\\[4pt]
\end{center}
\vspace{0.5cm}
\hrule height 1.5pt
\vspace{0.3cm}

% ============================================================
\section{Iterative Elimination of Weakly Dominated Strategies}

Recall from the previous lecture that the elimination of strictly dominated strategies one by one is called \textbf{Iterative Elimination of Strictly Dominated Strategies (IESDS)}. If this process gives a unique strategy profile, we call it a solution to the game.

We can apply a similar principle for \emph{weakly} dominated strategies.

\begin{warnbox}{Assumption 1'}
  A rational player does not use a \textbf{weakly dominated strategy}.
\end{warnbox}

\begin{defbox}{Rational Outcome}
  A strategy profile is called \textbf{rational} if it is the unique result of an iterative elimination of weakly dominated strategies.
\end{defbox}

\begin{notebox}
  As shown in the previous lecture, the order of elimination matters for weakly dominated strategies. If different orders lead to different unique outcomes, the solution might be ambiguous.
\end{notebox}

% ============================================================
\section{Second Price Sealed Bid Auction (Vickrey Auction)}

A classic application of weakly dominated strategies is the second-price sealed-bid auction.

\subsection{Formal Setup}
\begin{itemize}
    \item \textbf{Players:} $N = \{1, 2, \dots, n\}$
    \item \textbf{Strategies (Bids):} $S_i = [0, \infty)$ for each player $i$. Let player $i$'s bid be $b_i$.
    \item \textbf{Valuations:} Each player has a true valuation $v_i$ for the item.
    \item \textbf{Payoff Function:} $u_i(b_1, \dots, b_n)$ is defined as:
    \[
    u_i(b_1, \dots, b_n) = 
    \begin{cases}
      0 & \text{if } b_i < \max_{j \in N} b_j \\[8pt]
      \dfrac{v_i - \max_{j \neq i} b_j}{\left|\{k \mid b_k = \max_{j \in N} b_j\}\right|} & \text{if } b_i = \max_{j \in N} b_j
    \end{cases}
    \]
\end{itemize}

\begin{intubox}
  \textbf{How it works:}
  The highest bidder wins the item but pays the \emph{second-highest} bid ($\max_{j \neq i} b_j$). If there is a tie for the highest bid, the tied players split the surplus equally (represented by the denominator).
\end{intubox}

\begin{notebox}
  Since players do not know each other's valuations $v_j$, this is \textbf{not} a complete information game. It is a game of incomplete information.
\end{notebox}

\subsection{Truthful Bidding is a Weakly Dominant Strategy}

\begin{thmbox}{Vickrey's Theorem}
  In the second-price sealed-bid auction, \textbf{bidding your own valuation weakly dominates all other strategies.}
\end{thmbox}

In other words, for player $i$, playing $b_i = v_i$ is always at least as good as any other bid $b'_i$, regardless of what the other players do:
\[
  u_i(v_i, b_{-i}) \geq u_i(b'_i, b_{-i}) \quad \forall b_{-i} \in [0, \infty)^{n-1}, \; \forall b'_i \in [0, \infty)
\]

\textbf{Proof Reference:} See Section 4.6, Theorem 4.15 in Solan's Game Theory textbook for the detailed proof. The intuition is that changing your bid from $v_i$ only changes the outcome when it changes whether you win or lose, and in those cases, deviating from $v_i$ either makes you lose when winning was profitable, or makes you win when the price is higher than your valuation.

% ============================================================
\section{The Max-Min Approach}

Another way to approach a game is to think defensively. What is the absolute worst that can happen if I play a certain strategy, and how can I make that worst-case scenario as good as possible?

\begin{exbox}{Defensive Play Example}
  Consider the following game matrix from Player 1's perspective:
  \begin{center}
  \renewcommand{\arraystretch}{1.3}
  \begin{tabular}{c|cc|c}
    $P_1 \setminus P_2$ & L & R & \textbf{Row Min} \\
    \hline
    T & \payoff{2}{1} & \payoff{2}{-20} & \textbf{\color{warnred} 2} \\
    M & \payoff{3}{0} & \payoff{-10}{1} & -10 \\
    B & \payoff{-100}{2} & \payoff{3}{3} & -100 \\
  \end{tabular}
  \end{center}
  For each strategy, Player 1 looks at the minimum possible payoff:
  \begin{itemize}
      \item If T is played: min(2, 2) = 2
      \item If M is played: min(3, -10) = -10
      \item If B is played: min(-100, 3) = -100
  \end{itemize}
  The maximum of these minimums is 2, achieved by playing T. Strategy T guarantees Player 1 a payoff of at least 2, regardless of Player 2's action.
\end{exbox}

\subsection{Formal Definitions}

Suppose player $i$ plays strategy $s_i$. The minimum amount he/she will get is:
\[
  \min_{s_{-i} \in S_{-i}} u_i(s_i, s_{-i})
\]

\begin{defbox}{Max-Min Value}
  The \textbf{max-min value} of player $i$, denoted $\underline{v}_i$, is the highest payoff player $i$ can guarantee for themselves:
  \[
    \underline{v}_i = \max_{s_i \in S_i} \left[ \min_{s_{-i} \in S_{-i}} u_i(s_i, s_{-i}) \right]
  \]
\end{defbox}

\begin{defbox}{Max-Min Strategy}
  Suppose $s_i^* \in \arg\max_{s_i \in S_i} \left[ \min_{s_{-i} \in S_{-i}} u_i(s_i, s_{-i}) \right]$. 
  
  Then $s_i^*$ is called a \textbf{max-min strategy} for player $i$.
\end{defbox}

\begin{takebox}
  By playing a max-min strategy, a player can ensure they receive at least their max-min value, no matter how the opponents play:
  \[
    u_i(s_i^*, s_{-i}) \geq \underline{v}_i \quad \forall s_{-i} \in S_{-i}
  \]
\end{takebox}

\end{document}